% !TEX program = xelatex
\documentclass[a4paper,11pt,openany,fontset=none]{ctexrep}

\newcommand{\reporttitle}{收益目标策略：全产品尽职调查证据备忘录}
\newcommand{\reportsubtitle}{10 个已命名产品逐项穿透；研究框架通过，实施与配置未获批准}
\newcommand{\reportkicker}{CROSS-ASSET PRODUCT DUE DILIGENCE}
\newcommand{\reportscope}{中国交易所交易基金与公募 REIT | 10 个已命名产品}
\newcommand{\reportdate}{2026年8月12日}
\newcommand{\reportdatacutoff}{市场：2026-08-12；财务：2025 年报及 2026Q2}
\newcommand{\reporttype}{策略底层产品尽调 / 证据备忘录}
\newcommand{\reportauthor}{AStock Research}
\newcommand{\reporthouseview}{已完成 3 只 ETF 与 7 只公募 REIT 的一级资料产品卡、穿透暴露、历史现金与市场快照核验。结论不是“可以买什么”，而是：现有资料足以识别每一项真实经济暴露与关键缺口，但不足以证明 5\% 总回报或 3\% 净现金目标，也不足以放行任何市场交易产品。所有已命名产品维持 \textbf{q=0 / 仅观察}；两项未命名类别不可被当作已尽调标的。}
\newcommand{\reportquality}{一级资料：产品法律文件、2025 年报、2026Q2 报告与上交所同日市场快照；Q2 财务均未经审计。卖方资料已穷尽检索，但没有可用的产品级评级、目标或预测锚，估值权重为零。市场快照不含买卖价差、深度、IOPV/NAV 或容量。}
\newcommand{\reportdisclaimer}{本报告为公开资料研究与监控用途，不构成证券买卖建议、收益承诺或执行指令。}

% Reusable institutional-style preamble for XeLaTeX AStock research reports.
% Usage: define \reporttitle and \reportdate, then % Reusable institutional-style preamble for XeLaTeX AStock research reports.
% Usage: define \reporttitle and \reportdate, then % Reusable institutional-style preamble for XeLaTeX AStock research reports.
% Usage: define \reporttitle and \reportdate, then \input{.agents/templates/preamble.tex}.
% Design benchmark: JPM first-page dashboard, BlackRock house-view framing,
% Vanguard scenario discipline, and AB/JPM chartbook exhibit structure.

\providecommand{\reporttitle}{AStock Research Report}
\providecommand{\reportsubtitle}{House View, Evidence, Valuation and Risk}
\providecommand{\reportkicker}{INSTITUTIONAL EQUITY RESEARCH}
\providecommand{\reportscope}{CHINA A-SHARES}
\providecommand{\reportdate}{\today}
\providecommand{\reportdatacutoff}{Data cutoff not specified}
\providecommand{\reporttype}{Research note}
\providecommand{\reportauthor}{AStock Research Agent}
\providecommand{\reporthouseview}{State the house view in one decision-useful paragraph.}
\providecommand{\reportquality}{Source quality not specified.}
\providecommand{\reportdisclaimer}{本报告基于公开资料整理，不构成任何证券买卖建议。}

% ============ 页面布局 ============
\usepackage[top=1.8cm, bottom=1.9cm, left=1.8cm, right=1.8cm]{geometry}
\setlength{\parskip}{0.3em}
\setlength{\parindent}{2em}
\renewcommand{\baselinestretch}{1.12}
\raggedbottom

% ============ Fonts ============
% Prefer the installed Noto CJK family when available (for Linux builds).
% Otherwise retain the original MacTeX typography. This avoids platform
% probing via shell escape and keeps the source portable across both engines.
\IfFontExistsTF{Noto Serif CJK SC}{
  \setCJKmainfont{Noto Serif CJK SC}
  \setCJKsansfont{Noto Sans CJK SC}
  \setCJKmonofont{Noto Sans CJK SC}
  \setmainfont{Liberation Serif}
  \newcommand{\kaishu}{\CJKfontspec{Noto Serif CJK SC}}
  \newcommand{\heiti}{\CJKfontspec{Noto Sans CJK SC}}
  \newcommand{\songti}{\CJKfontspec{Noto Serif CJK SC}}
}{
  \setCJKmainfont[BoldFont={Heiti SC}]{STSong}
  \setCJKsansfont{Heiti SC}
  \setCJKmonofont{Heiti SC}
  \setmainfont{Times New Roman}
  \newcommand{\kaishu}{\CJKfontspec{STKaiti}}
  \newcommand{\heiti}{\CJKfontspec{Heiti SC}}
  \newcommand{\songti}{\CJKfontspec{STSong}}
}

% ============ 颜色（IB Navy/Grey palette） ============
\usepackage{xcolor}
\definecolor{navy}{HTML}{003366}
\definecolor{deepnavy}{HTML}{0B1F33}
\definecolor{steelgrey}{HTML}{4A5568}
\definecolor{mutedgrey}{HTML}{718096}
\definecolor{lightgrey}{HTML}{F7FAFC}
\definecolor{panelgrey}{HTML}{EDF2F7}
\definecolor{riskred}{HTML}{C53030}
\definecolor{riskamber}{HTML}{D69E2E}
\definecolor{riskgreen}{HTML}{38A169}
\definecolor{riskblue}{HTML}{2B6CB0}
\definecolor{accentblue}{HTML}{2B6CB0}
\definecolor{nvgreen}{HTML}{76B900}
\definecolor{rulegrey}{HTML}{CBD5E0}

% ============ 表格 ============
\usepackage{booktabs}
\usepackage{longtable}
\usepackage{tabularx}
\usepackage{multirow}
\usepackage{array}
\usepackage{makecell}
\usepackage{colortbl}
\renewcommand{\arraystretch}{1.12}
\newcolumntype{Y}{>{\centering\arraybackslash}X}
\newcolumntype{L}[1]{>{\raggedright\arraybackslash}p{#1}}
\newcolumntype{C}[1]{>{\centering\arraybackslash}p{#1}}
\newcolumntype{R}[1]{>{\raggedleft\arraybackslash}p{#1}}

% ============ 图形 ============
\usepackage{graphicx}
\usepackage{tikz}
\usetikzlibrary{shapes,arrows.meta,positioning,calc,backgrounds,fit,decorations.pathreplacing}
\usepackage{pgfplots}
\usepgfplotslibrary{polar}
\pgfplotsset{compat=1.18}
\usepackage[edges]{forest}

% ============ 页眉页脚（报告标题和日期需在main.tex中覆盖） ============
\usepackage{fancyhdr}
\pagestyle{fancy}
\fancyhf{}
\fancyhead[L]{\scriptsize\color{steelgrey} \reporttitle}
\fancyhead[R]{\scriptsize\color{steelgrey} \reportdate}
\fancyfoot[C]{\small\color{steelgrey} \thepage}
\fancyfoot[R]{\tiny\color{steelgrey} \reportdisclaimer}
\renewcommand{\headrulewidth}{0.4pt}
\renewcommand{\footrulewidth}{0.2pt}

% ============ 章节格式 ============
\usepackage{titlesec}
\titleformat{\chapter}[block]
  {\normalfont\LARGE\bfseries\color{navy}}
  {\thechapter}{0.9em}{}
\titlespacing*{\chapter}{0pt}{0.4em}{1.0em}
\titleformat{\section}
  {\normalfont\large\bfseries\color{navy}}
  {\thesection}{1em}{}
\titleformat{\subsection}
  {\normalfont\normalsize\bfseries\color{steelgrey}}
  {\thesubsection}{1em}{}

% ============ 超链接 ============
\usepackage{hyperref}
\hypersetup{
  colorlinks=true,
  linkcolor=navy,
  urlcolor=navy,
  citecolor=navy,
}

% ============ 其他工具 ============
\usepackage{enumitem}
\setlist{nosep, leftmargin=2em}
\IfFileExists{fontawesome5.sty}{\usepackage{fontawesome5}}{}
% ============ Class B 全局兜底（xeCJK + 缺字 fallback 阻断）：
% 当 fontawesome5 因 Times New Roman 缺字 / \FA 命名空间失败时，
% 禁止输出 aExclamationTriangle / aCloud 裸文本，一律降级为 amssymb 原生符号。
% 注：amssymb 在 L112 已加载，此处不重复。所有 \def 在本段落。
\def\faExclamationTriangle{\ensuremath{\textcolor{riskred}{\blacktriangle}}\kern-0.25em}
\def\faCloud{\ensuremath{\clubsuit}\kern-0.1em}
\def\faYenSign{\ensuremath{\yen}}
\def\faIndustry{\ensuremath{\S}}
\def\faCogs{\ensuremath{\ast}}
\def\faFlagCheckered{\ensuremath{\blacktriangleright}}
\def\faMemory{\textbf{[RAM]}}
\def\faServer{\textbf{[DC]}}
\def\faGlobe{\ensuremath{\diamond}}
\def\faStar{\ensuremath{\bigstar}}
\def\faMicrochip{\textbf{[$\mu$P]}}
\def\faLayerGroup{\textbf{[LAY]}}
\def\faCubes{\ensuremath{\boxtimes}}
\usepackage{amssymb}
\usepackage{pifont}
\usepackage{tcolorbox}
\tcbuselibrary{skins,breakable}
\usepackage{caption}
\captionsetup{font=small, skip=4pt}
\usepackage{float}
\usepackage{needspace}
\setlength{\textfloatsep}{8pt plus 2pt minus 2pt}
\setlength{\floatsep}{6pt plus 2pt minus 2pt}
\setlength{\intextsep}{6pt plus 2pt minus 2pt}
\setlength{\abovecaptionskip}{4pt}
\setlength{\belowcaptionskip}{2pt}

% ============ 自定义命令 ============
\newcommand{\rmark}{\textcolor{riskred}{\ding{108}}}
\newcommand{\amark}{\textcolor{riskamber}{\ding{108}}}
\newcommand{\gmark}{\textcolor{riskgreen}{\ding{108}}}
\newcommand{\starfive}{$\bigstar\bigstar\bigstar\bigstar\bigstar$}
\newcommand{\starfour}{$\bigstar\bigstar\bigstar\bigstar$}
\newcommand{\starthree}{$\bigstar\bigstar\bigstar$}
% 资料来源注：使用楷体（STKaiti）而非斜体，避免中文无斜体字形导致的字体回退警告
% \mdseries 避免继承周围 \bfseries 导致 STKaiti 粗体缺失警告
% \sloppy + \hbadness 允许长文件名/URL 在必要时断行
% 局部重定义 \textbf → \heiti：STKaiti 无粗体，用黑体替代强调
\newcommand{\sourcenote}[1]{\par{\scriptsize\mdseries\color{mutedgrey}\kaishu\sloppy\hbadness=9999\def\textbf##1{{\heiti ##1}} 资料来源：#1}\vspace{0.2em}}
\newcommand{\qualitynote}[1]{\par{\scriptsize\mdseries\color{mutedgrey}\sloppy\def\textbf##1{{\heiti ##1}} 数据质量：#1}\vspace{0.2em}}
\newcommand{\evidencetag}[1]{%
  \tikz[baseline=(x.base)]\node[
    xshift=0pt,
    inner xsep=3.5pt,
    inner ysep=1.2pt,
    rounded corners=1.5pt,
    fill=panelgrey,
    draw=rulegrey,
    text=steelgrey,
    font=\scriptsize
  ] (x) {#1};%
}
\newcommand{\stance}[2]{\textbf{\textcolor{#1}{#2}}}

% 信息框
\newtcolorbox{keyinsight}[1][]{
  colback=lightgrey, colframe=navy,
  fonttitle=\bfseries, title={#1}, breakable,
  arc=2pt, boxrule=0.5pt, left=8pt, right=8pt, top=6pt, bottom=6pt
}
\newtcolorbox{riskbox}[1][]{
  colback=red!5, colframe=riskred,
  fonttitle=\bfseries, title={#1}, breakable,
  arc=2pt, boxrule=0.5pt, left=8pt, right=8pt, top=6pt, bottom=6pt
}
\newtcolorbox{dashboardbox}[1][]{
  enhanced, breakable,
  colback=white, colframe=navy,
  boxrule=0.7pt, arc=2pt,
  fonttitle=\bfseries\color{white},
  coltitle=white, colbacktitle=navy,
  title={#1},
  left=9pt, right=9pt, top=7pt, bottom=7pt
}
\newtcolorbox{houseviewbox}[1][]{
  enhanced, breakable,
  colback=lightgrey, colframe=deepnavy,
  borderline west={3pt}{0pt}{accentblue},
  boxrule=0.4pt, arc=2pt,
  fonttitle=\bfseries, title={#1},
  left=9pt, right=9pt, top=7pt, bottom=7pt
}
\newtcolorbox{sourcequalitybox}[1][]{
  enhanced, breakable,
  colback=yellow!6, colframe=riskamber,
  boxrule=0.45pt, arc=2pt,
  fonttitle=\bfseries, title={#1},
  left=9pt, right=9pt, top=7pt, bottom=7pt
}
% exhibitbox：表/图外框。left=7pt+right=7pt=14pt 内边距，before upper 内的 \linewidth 已扣除该内边距。
% 约定：\exhibitboxwidth = 真实内边距宽（tcolorbox 已扣好 padding 的 \linewidth）。
% 由于 tabularx/longtable 宽度参数的语义是「外部总宽（含两侧 \tabcolsep）」，内部列实际可用 = W - 2\tabcolsep。
% 所以 tabularx 直接传 \exhibitboxwidth（总宽 == 内边距宽），而其他用途（resizebox、tikz、p{...}）用 \exhibitwidth = \exhibitboxwidth - 2\tabcolsep。
\newlength{\exhibitboxwidth}
\newlength{\exhibitwidth}
\newtcolorbox{exhibitbox}[1][]{
  enhanced, breakable,
  colback=white, colframe=rulegrey,
  borderline north={1.2pt}{0pt}{navy},
  boxrule=0.35pt, arc=1pt,
  fonttitle=\bfseries\color{navy}, title={#1},
  left=7pt, right=7pt, top=6pt, bottom=6pt,
  before upper={\global\exhibitboxwidth=\linewidth\relax
    \global\exhibitwidth=\dimexpr\linewidth-2\tabcolsep\relax},
}
\newtcolorbox{disclosurebox}[1][]{
  enhanced, breakable,
  colback=lightgrey, colframe=steelgrey,
  boxrule=0.35pt, arc=1pt,
  fonttitle=\bfseries, title={#1},
  left=8pt, right=8pt, top=6pt, bottom=6pt
}

% ============ First-page dashboard cover ============
\newcommand{\astockcover}{%
\begin{titlepage}
\thispagestyle{empty}
\begin{tikzpicture}[remember picture, overlay]
  \fill[deepnavy] (current page.north west) rectangle ([yshift=-2.55cm]current page.north east);
  \fill[accentblue] ([yshift=-2.55cm]current page.north west) rectangle ([yshift=-2.66cm]current page.north east);
  \fill[deepnavy] (current page.south west) rectangle ([yshift=1.2cm]current page.south east);
\end{tikzpicture}
\vspace*{0.62cm}
{\color{white}\sffamily\bfseries\small \reportkicker}\\[0.18cm]
{\color{steelgrey}\sffamily\scriptsize \reportscope \quad|\quad \reporttype \quad|\quad \reportdate}\\[1.0cm]
{\fontsize{25}{30}\selectfont\bfseries\color{navy}\reporttitle}\\[0.25cm]
{\fontsize{13}{17}\selectfont\color{steelgrey}\reportsubtitle}\\[0.75cm]
\begin{dashboardbox}[决策摘要]
\begin{tabularx}{\textwidth}{L{3.0cm}X L{3.0cm}X}
\textbf{报告类型} & \reporttype & \textbf{数据截至} & \reportdatacutoff \\
\textbf{覆盖范围} & \reportscope & \textbf{作者} & \reportauthor \\
\end{tabularx}
\vspace{0.45em}
\begin{houseviewbox}[本机构观点]
\reporthouseview
\end{houseviewbox}
\begin{sourcequalitybox}[证据质量]
\reportquality
\end{sourcequalitybox}
\end{dashboardbox}
\vfill
{\scriptsize\color{steelgrey}\reportdisclaimer}
\end{titlepage}
}
.
% Design benchmark: JPM first-page dashboard, BlackRock house-view framing,
% Vanguard scenario discipline, and AB/JPM chartbook exhibit structure.

\providecommand{\reporttitle}{AStock Research Report}
\providecommand{\reportsubtitle}{House View, Evidence, Valuation and Risk}
\providecommand{\reportkicker}{INSTITUTIONAL EQUITY RESEARCH}
\providecommand{\reportscope}{CHINA A-SHARES}
\providecommand{\reportdate}{\today}
\providecommand{\reportdatacutoff}{Data cutoff not specified}
\providecommand{\reporttype}{Research note}
\providecommand{\reportauthor}{AStock Research Agent}
\providecommand{\reporthouseview}{State the house view in one decision-useful paragraph.}
\providecommand{\reportquality}{Source quality not specified.}
\providecommand{\reportdisclaimer}{本报告基于公开资料整理，不构成任何证券买卖建议。}

% ============ 页面布局 ============
\usepackage[top=1.8cm, bottom=1.9cm, left=1.8cm, right=1.8cm]{geometry}
\setlength{\parskip}{0.3em}
\setlength{\parindent}{2em}
\renewcommand{\baselinestretch}{1.12}
\raggedbottom

% ============ Fonts ============
% Prefer the installed Noto CJK family when available (for Linux builds).
% Otherwise retain the original MacTeX typography. This avoids platform
% probing via shell escape and keeps the source portable across both engines.
\IfFontExistsTF{Noto Serif CJK SC}{
  \setCJKmainfont{Noto Serif CJK SC}
  \setCJKsansfont{Noto Sans CJK SC}
  \setCJKmonofont{Noto Sans CJK SC}
  \setmainfont{Liberation Serif}
  \newcommand{\kaishu}{\CJKfontspec{Noto Serif CJK SC}}
  \newcommand{\heiti}{\CJKfontspec{Noto Sans CJK SC}}
  \newcommand{\songti}{\CJKfontspec{Noto Serif CJK SC}}
}{
  \setCJKmainfont[BoldFont={Heiti SC}]{STSong}
  \setCJKsansfont{Heiti SC}
  \setCJKmonofont{Heiti SC}
  \setmainfont{Times New Roman}
  \newcommand{\kaishu}{\CJKfontspec{STKaiti}}
  \newcommand{\heiti}{\CJKfontspec{Heiti SC}}
  \newcommand{\songti}{\CJKfontspec{STSong}}
}

% ============ 颜色（IB Navy/Grey palette） ============
\usepackage{xcolor}
\definecolor{navy}{HTML}{003366}
\definecolor{deepnavy}{HTML}{0B1F33}
\definecolor{steelgrey}{HTML}{4A5568}
\definecolor{mutedgrey}{HTML}{718096}
\definecolor{lightgrey}{HTML}{F7FAFC}
\definecolor{panelgrey}{HTML}{EDF2F7}
\definecolor{riskred}{HTML}{C53030}
\definecolor{riskamber}{HTML}{D69E2E}
\definecolor{riskgreen}{HTML}{38A169}
\definecolor{riskblue}{HTML}{2B6CB0}
\definecolor{accentblue}{HTML}{2B6CB0}
\definecolor{nvgreen}{HTML}{76B900}
\definecolor{rulegrey}{HTML}{CBD5E0}

% ============ 表格 ============
\usepackage{booktabs}
\usepackage{longtable}
\usepackage{tabularx}
\usepackage{multirow}
\usepackage{array}
\usepackage{makecell}
\usepackage{colortbl}
\renewcommand{\arraystretch}{1.12}
\newcolumntype{Y}{>{\centering\arraybackslash}X}
\newcolumntype{L}[1]{>{\raggedright\arraybackslash}p{#1}}
\newcolumntype{C}[1]{>{\centering\arraybackslash}p{#1}}
\newcolumntype{R}[1]{>{\raggedleft\arraybackslash}p{#1}}

% ============ 图形 ============
\usepackage{graphicx}
\usepackage{tikz}
\usetikzlibrary{shapes,arrows.meta,positioning,calc,backgrounds,fit,decorations.pathreplacing}
\usepackage{pgfplots}
\usepgfplotslibrary{polar}
\pgfplotsset{compat=1.18}
\usepackage[edges]{forest}

% ============ 页眉页脚（报告标题和日期需在main.tex中覆盖） ============
\usepackage{fancyhdr}
\pagestyle{fancy}
\fancyhf{}
\fancyhead[L]{\scriptsize\color{steelgrey} \reporttitle}
\fancyhead[R]{\scriptsize\color{steelgrey} \reportdate}
\fancyfoot[C]{\small\color{steelgrey} \thepage}
\fancyfoot[R]{\tiny\color{steelgrey} \reportdisclaimer}
\renewcommand{\headrulewidth}{0.4pt}
\renewcommand{\footrulewidth}{0.2pt}

% ============ 章节格式 ============
\usepackage{titlesec}
\titleformat{\chapter}[block]
  {\normalfont\LARGE\bfseries\color{navy}}
  {\thechapter}{0.9em}{}
\titlespacing*{\chapter}{0pt}{0.4em}{1.0em}
\titleformat{\section}
  {\normalfont\large\bfseries\color{navy}}
  {\thesection}{1em}{}
\titleformat{\subsection}
  {\normalfont\normalsize\bfseries\color{steelgrey}}
  {\thesubsection}{1em}{}

% ============ 超链接 ============
\usepackage{hyperref}
\hypersetup{
  colorlinks=true,
  linkcolor=navy,
  urlcolor=navy,
  citecolor=navy,
}

% ============ 其他工具 ============
\usepackage{enumitem}
\setlist{nosep, leftmargin=2em}
\IfFileExists{fontawesome5.sty}{\usepackage{fontawesome5}}{}
% ============ Class B 全局兜底（xeCJK + 缺字 fallback 阻断）：
% 当 fontawesome5 因 Times New Roman 缺字 / \FA 命名空间失败时，
% 禁止输出 aExclamationTriangle / aCloud 裸文本，一律降级为 amssymb 原生符号。
% 注：amssymb 在 L112 已加载，此处不重复。所有 \def 在本段落。
\def\faExclamationTriangle{\ensuremath{\textcolor{riskred}{\blacktriangle}}\kern-0.25em}
\def\faCloud{\ensuremath{\clubsuit}\kern-0.1em}
\def\faYenSign{\ensuremath{\yen}}
\def\faIndustry{\ensuremath{\S}}
\def\faCogs{\ensuremath{\ast}}
\def\faFlagCheckered{\ensuremath{\blacktriangleright}}
\def\faMemory{\textbf{[RAM]}}
\def\faServer{\textbf{[DC]}}
\def\faGlobe{\ensuremath{\diamond}}
\def\faStar{\ensuremath{\bigstar}}
\def\faMicrochip{\textbf{[$\mu$P]}}
\def\faLayerGroup{\textbf{[LAY]}}
\def\faCubes{\ensuremath{\boxtimes}}
\usepackage{amssymb}
\usepackage{pifont}
\usepackage{tcolorbox}
\tcbuselibrary{skins,breakable}
\usepackage{caption}
\captionsetup{font=small, skip=4pt}
\usepackage{float}
\usepackage{needspace}
\setlength{\textfloatsep}{8pt plus 2pt minus 2pt}
\setlength{\floatsep}{6pt plus 2pt minus 2pt}
\setlength{\intextsep}{6pt plus 2pt minus 2pt}
\setlength{\abovecaptionskip}{4pt}
\setlength{\belowcaptionskip}{2pt}

% ============ 自定义命令 ============
\newcommand{\rmark}{\textcolor{riskred}{\ding{108}}}
\newcommand{\amark}{\textcolor{riskamber}{\ding{108}}}
\newcommand{\gmark}{\textcolor{riskgreen}{\ding{108}}}
\newcommand{\starfive}{$\bigstar\bigstar\bigstar\bigstar\bigstar$}
\newcommand{\starfour}{$\bigstar\bigstar\bigstar\bigstar$}
\newcommand{\starthree}{$\bigstar\bigstar\bigstar$}
% 资料来源注：使用楷体（STKaiti）而非斜体，避免中文无斜体字形导致的字体回退警告
% \mdseries 避免继承周围 \bfseries 导致 STKaiti 粗体缺失警告
% \sloppy + \hbadness 允许长文件名/URL 在必要时断行
% 局部重定义 \textbf → \heiti：STKaiti 无粗体，用黑体替代强调
\newcommand{\sourcenote}[1]{\par{\scriptsize\mdseries\color{mutedgrey}\kaishu\sloppy\hbadness=9999\def\textbf##1{{\heiti ##1}} 资料来源：#1}\vspace{0.2em}}
\newcommand{\qualitynote}[1]{\par{\scriptsize\mdseries\color{mutedgrey}\sloppy\def\textbf##1{{\heiti ##1}} 数据质量：#1}\vspace{0.2em}}
\newcommand{\evidencetag}[1]{%
  \tikz[baseline=(x.base)]\node[
    xshift=0pt,
    inner xsep=3.5pt,
    inner ysep=1.2pt,
    rounded corners=1.5pt,
    fill=panelgrey,
    draw=rulegrey,
    text=steelgrey,
    font=\scriptsize
  ] (x) {#1};%
}
\newcommand{\stance}[2]{\textbf{\textcolor{#1}{#2}}}

% 信息框
\newtcolorbox{keyinsight}[1][]{
  colback=lightgrey, colframe=navy,
  fonttitle=\bfseries, title={#1}, breakable,
  arc=2pt, boxrule=0.5pt, left=8pt, right=8pt, top=6pt, bottom=6pt
}
\newtcolorbox{riskbox}[1][]{
  colback=red!5, colframe=riskred,
  fonttitle=\bfseries, title={#1}, breakable,
  arc=2pt, boxrule=0.5pt, left=8pt, right=8pt, top=6pt, bottom=6pt
}
\newtcolorbox{dashboardbox}[1][]{
  enhanced, breakable,
  colback=white, colframe=navy,
  boxrule=0.7pt, arc=2pt,
  fonttitle=\bfseries\color{white},
  coltitle=white, colbacktitle=navy,
  title={#1},
  left=9pt, right=9pt, top=7pt, bottom=7pt
}
\newtcolorbox{houseviewbox}[1][]{
  enhanced, breakable,
  colback=lightgrey, colframe=deepnavy,
  borderline west={3pt}{0pt}{accentblue},
  boxrule=0.4pt, arc=2pt,
  fonttitle=\bfseries, title={#1},
  left=9pt, right=9pt, top=7pt, bottom=7pt
}
\newtcolorbox{sourcequalitybox}[1][]{
  enhanced, breakable,
  colback=yellow!6, colframe=riskamber,
  boxrule=0.45pt, arc=2pt,
  fonttitle=\bfseries, title={#1},
  left=9pt, right=9pt, top=7pt, bottom=7pt
}
% exhibitbox：表/图外框。left=7pt+right=7pt=14pt 内边距，before upper 内的 \linewidth 已扣除该内边距。
% 约定：\exhibitboxwidth = 真实内边距宽（tcolorbox 已扣好 padding 的 \linewidth）。
% 由于 tabularx/longtable 宽度参数的语义是「外部总宽（含两侧 \tabcolsep）」，内部列实际可用 = W - 2\tabcolsep。
% 所以 tabularx 直接传 \exhibitboxwidth（总宽 == 内边距宽），而其他用途（resizebox、tikz、p{...}）用 \exhibitwidth = \exhibitboxwidth - 2\tabcolsep。
\newlength{\exhibitboxwidth}
\newlength{\exhibitwidth}
\newtcolorbox{exhibitbox}[1][]{
  enhanced, breakable,
  colback=white, colframe=rulegrey,
  borderline north={1.2pt}{0pt}{navy},
  boxrule=0.35pt, arc=1pt,
  fonttitle=\bfseries\color{navy}, title={#1},
  left=7pt, right=7pt, top=6pt, bottom=6pt,
  before upper={\global\exhibitboxwidth=\linewidth\relax
    \global\exhibitwidth=\dimexpr\linewidth-2\tabcolsep\relax},
}
\newtcolorbox{disclosurebox}[1][]{
  enhanced, breakable,
  colback=lightgrey, colframe=steelgrey,
  boxrule=0.35pt, arc=1pt,
  fonttitle=\bfseries, title={#1},
  left=8pt, right=8pt, top=6pt, bottom=6pt
}

% ============ First-page dashboard cover ============
\newcommand{\astockcover}{%
\begin{titlepage}
\thispagestyle{empty}
\begin{tikzpicture}[remember picture, overlay]
  \fill[deepnavy] (current page.north west) rectangle ([yshift=-2.55cm]current page.north east);
  \fill[accentblue] ([yshift=-2.55cm]current page.north west) rectangle ([yshift=-2.66cm]current page.north east);
  \fill[deepnavy] (current page.south west) rectangle ([yshift=1.2cm]current page.south east);
\end{tikzpicture}
\vspace*{0.62cm}
{\color{white}\sffamily\bfseries\small \reportkicker}\\[0.18cm]
{\color{steelgrey}\sffamily\scriptsize \reportscope \quad|\quad \reporttype \quad|\quad \reportdate}\\[1.0cm]
{\fontsize{25}{30}\selectfont\bfseries\color{navy}\reporttitle}\\[0.25cm]
{\fontsize{13}{17}\selectfont\color{steelgrey}\reportsubtitle}\\[0.75cm]
\begin{dashboardbox}[决策摘要]
\begin{tabularx}{\textwidth}{L{3.0cm}X L{3.0cm}X}
\textbf{报告类型} & \reporttype & \textbf{数据截至} & \reportdatacutoff \\
\textbf{覆盖范围} & \reportscope & \textbf{作者} & \reportauthor \\
\end{tabularx}
\vspace{0.45em}
\begin{houseviewbox}[本机构观点]
\reporthouseview
\end{houseviewbox}
\begin{sourcequalitybox}[证据质量]
\reportquality
\end{sourcequalitybox}
\end{dashboardbox}
\vfill
{\scriptsize\color{steelgrey}\reportdisclaimer}
\end{titlepage}
}
.
% Design benchmark: JPM first-page dashboard, BlackRock house-view framing,
% Vanguard scenario discipline, and AB/JPM chartbook exhibit structure.

\providecommand{\reporttitle}{AStock Research Report}
\providecommand{\reportsubtitle}{House View, Evidence, Valuation and Risk}
\providecommand{\reportkicker}{INSTITUTIONAL EQUITY RESEARCH}
\providecommand{\reportscope}{CHINA A-SHARES}
\providecommand{\reportdate}{\today}
\providecommand{\reportdatacutoff}{Data cutoff not specified}
\providecommand{\reporttype}{Research note}
\providecommand{\reportauthor}{AStock Research Agent}
\providecommand{\reporthouseview}{State the house view in one decision-useful paragraph.}
\providecommand{\reportquality}{Source quality not specified.}
\providecommand{\reportdisclaimer}{本报告基于公开资料整理，不构成任何证券买卖建议。}

% ============ 页面布局 ============
\usepackage[top=1.8cm, bottom=1.9cm, left=1.8cm, right=1.8cm]{geometry}
\setlength{\parskip}{0.3em}
\setlength{\parindent}{2em}
\renewcommand{\baselinestretch}{1.12}
\raggedbottom

% ============ Fonts ============
% Prefer the installed Noto CJK family when available (for Linux builds).
% Otherwise retain the original MacTeX typography. This avoids platform
% probing via shell escape and keeps the source portable across both engines.
\IfFontExistsTF{Noto Serif CJK SC}{
  \setCJKmainfont{Noto Serif CJK SC}
  \setCJKsansfont{Noto Sans CJK SC}
  \setCJKmonofont{Noto Sans CJK SC}
  \setmainfont{Liberation Serif}
  \newcommand{\kaishu}{\CJKfontspec{Noto Serif CJK SC}}
  \newcommand{\heiti}{\CJKfontspec{Noto Sans CJK SC}}
  \newcommand{\songti}{\CJKfontspec{Noto Serif CJK SC}}
}{
  \setCJKmainfont[BoldFont={Heiti SC}]{STSong}
  \setCJKsansfont{Heiti SC}
  \setCJKmonofont{Heiti SC}
  \setmainfont{Times New Roman}
  \newcommand{\kaishu}{\CJKfontspec{STKaiti}}
  \newcommand{\heiti}{\CJKfontspec{Heiti SC}}
  \newcommand{\songti}{\CJKfontspec{STSong}}
}

% ============ 颜色（IB Navy/Grey palette） ============
\usepackage{xcolor}
\definecolor{navy}{HTML}{003366}
\definecolor{deepnavy}{HTML}{0B1F33}
\definecolor{steelgrey}{HTML}{4A5568}
\definecolor{mutedgrey}{HTML}{718096}
\definecolor{lightgrey}{HTML}{F7FAFC}
\definecolor{panelgrey}{HTML}{EDF2F7}
\definecolor{riskred}{HTML}{C53030}
\definecolor{riskamber}{HTML}{D69E2E}
\definecolor{riskgreen}{HTML}{38A169}
\definecolor{riskblue}{HTML}{2B6CB0}
\definecolor{accentblue}{HTML}{2B6CB0}
\definecolor{nvgreen}{HTML}{76B900}
\definecolor{rulegrey}{HTML}{CBD5E0}

% ============ 表格 ============
\usepackage{booktabs}
\usepackage{longtable}
\usepackage{tabularx}
\usepackage{multirow}
\usepackage{array}
\usepackage{makecell}
\usepackage{colortbl}
\renewcommand{\arraystretch}{1.12}
\newcolumntype{Y}{>{\centering\arraybackslash}X}
\newcolumntype{L}[1]{>{\raggedright\arraybackslash}p{#1}}
\newcolumntype{C}[1]{>{\centering\arraybackslash}p{#1}}
\newcolumntype{R}[1]{>{\raggedleft\arraybackslash}p{#1}}

% ============ 图形 ============
\usepackage{graphicx}
\usepackage{tikz}
\usetikzlibrary{shapes,arrows.meta,positioning,calc,backgrounds,fit,decorations.pathreplacing}
\usepackage{pgfplots}
\usepgfplotslibrary{polar}
\pgfplotsset{compat=1.18}
\usepackage[edges]{forest}

% ============ 页眉页脚（报告标题和日期需在main.tex中覆盖） ============
\usepackage{fancyhdr}
\pagestyle{fancy}
\fancyhf{}
\fancyhead[L]{\scriptsize\color{steelgrey} \reporttitle}
\fancyhead[R]{\scriptsize\color{steelgrey} \reportdate}
\fancyfoot[C]{\small\color{steelgrey} \thepage}
\fancyfoot[R]{\tiny\color{steelgrey} \reportdisclaimer}
\renewcommand{\headrulewidth}{0.4pt}
\renewcommand{\footrulewidth}{0.2pt}

% ============ 章节格式 ============
\usepackage{titlesec}
\titleformat{\chapter}[block]
  {\normalfont\LARGE\bfseries\color{navy}}
  {\thechapter}{0.9em}{}
\titlespacing*{\chapter}{0pt}{0.4em}{1.0em}
\titleformat{\section}
  {\normalfont\large\bfseries\color{navy}}
  {\thesection}{1em}{}
\titleformat{\subsection}
  {\normalfont\normalsize\bfseries\color{steelgrey}}
  {\thesubsection}{1em}{}

% ============ 超链接 ============
\usepackage{hyperref}
\hypersetup{
  colorlinks=true,
  linkcolor=navy,
  urlcolor=navy,
  citecolor=navy,
}

% ============ 其他工具 ============
\usepackage{enumitem}
\setlist{nosep, leftmargin=2em}
\IfFileExists{fontawesome5.sty}{\usepackage{fontawesome5}}{}
% ============ Class B 全局兜底（xeCJK + 缺字 fallback 阻断）：
% 当 fontawesome5 因 Times New Roman 缺字 / \FA 命名空间失败时，
% 禁止输出 aExclamationTriangle / aCloud 裸文本，一律降级为 amssymb 原生符号。
% 注：amssymb 在 L112 已加载，此处不重复。所有 \def 在本段落。
\def\faExclamationTriangle{\ensuremath{\textcolor{riskred}{\blacktriangle}}\kern-0.25em}
\def\faCloud{\ensuremath{\clubsuit}\kern-0.1em}
\def\faYenSign{\ensuremath{\yen}}
\def\faIndustry{\ensuremath{\S}}
\def\faCogs{\ensuremath{\ast}}
\def\faFlagCheckered{\ensuremath{\blacktriangleright}}
\def\faMemory{\textbf{[RAM]}}
\def\faServer{\textbf{[DC]}}
\def\faGlobe{\ensuremath{\diamond}}
\def\faStar{\ensuremath{\bigstar}}
\def\faMicrochip{\textbf{[$\mu$P]}}
\def\faLayerGroup{\textbf{[LAY]}}
\def\faCubes{\ensuremath{\boxtimes}}
\usepackage{amssymb}
\usepackage{pifont}
\usepackage{tcolorbox}
\tcbuselibrary{skins,breakable}
\usepackage{caption}
\captionsetup{font=small, skip=4pt}
\usepackage{float}
\usepackage{needspace}
\setlength{\textfloatsep}{8pt plus 2pt minus 2pt}
\setlength{\floatsep}{6pt plus 2pt minus 2pt}
\setlength{\intextsep}{6pt plus 2pt minus 2pt}
\setlength{\abovecaptionskip}{4pt}
\setlength{\belowcaptionskip}{2pt}

% ============ 自定义命令 ============
\newcommand{\rmark}{\textcolor{riskred}{\ding{108}}}
\newcommand{\amark}{\textcolor{riskamber}{\ding{108}}}
\newcommand{\gmark}{\textcolor{riskgreen}{\ding{108}}}
\newcommand{\starfive}{$\bigstar\bigstar\bigstar\bigstar\bigstar$}
\newcommand{\starfour}{$\bigstar\bigstar\bigstar\bigstar$}
\newcommand{\starthree}{$\bigstar\bigstar\bigstar$}
% 资料来源注：使用楷体（STKaiti）而非斜体，避免中文无斜体字形导致的字体回退警告
% \mdseries 避免继承周围 \bfseries 导致 STKaiti 粗体缺失警告
% \sloppy + \hbadness 允许长文件名/URL 在必要时断行
% 局部重定义 \textbf → \heiti：STKaiti 无粗体，用黑体替代强调
\newcommand{\sourcenote}[1]{\par{\scriptsize\mdseries\color{mutedgrey}\kaishu\sloppy\hbadness=9999\def\textbf##1{{\heiti ##1}} 资料来源：#1}\vspace{0.2em}}
\newcommand{\qualitynote}[1]{\par{\scriptsize\mdseries\color{mutedgrey}\sloppy\def\textbf##1{{\heiti ##1}} 数据质量：#1}\vspace{0.2em}}
\newcommand{\evidencetag}[1]{%
  \tikz[baseline=(x.base)]\node[
    xshift=0pt,
    inner xsep=3.5pt,
    inner ysep=1.2pt,
    rounded corners=1.5pt,
    fill=panelgrey,
    draw=rulegrey,
    text=steelgrey,
    font=\scriptsize
  ] (x) {#1};%
}
\newcommand{\stance}[2]{\textbf{\textcolor{#1}{#2}}}

% 信息框
\newtcolorbox{keyinsight}[1][]{
  colback=lightgrey, colframe=navy,
  fonttitle=\bfseries, title={#1}, breakable,
  arc=2pt, boxrule=0.5pt, left=8pt, right=8pt, top=6pt, bottom=6pt
}
\newtcolorbox{riskbox}[1][]{
  colback=red!5, colframe=riskred,
  fonttitle=\bfseries, title={#1}, breakable,
  arc=2pt, boxrule=0.5pt, left=8pt, right=8pt, top=6pt, bottom=6pt
}
\newtcolorbox{dashboardbox}[1][]{
  enhanced, breakable,
  colback=white, colframe=navy,
  boxrule=0.7pt, arc=2pt,
  fonttitle=\bfseries\color{white},
  coltitle=white, colbacktitle=navy,
  title={#1},
  left=9pt, right=9pt, top=7pt, bottom=7pt
}
\newtcolorbox{houseviewbox}[1][]{
  enhanced, breakable,
  colback=lightgrey, colframe=deepnavy,
  borderline west={3pt}{0pt}{accentblue},
  boxrule=0.4pt, arc=2pt,
  fonttitle=\bfseries, title={#1},
  left=9pt, right=9pt, top=7pt, bottom=7pt
}
\newtcolorbox{sourcequalitybox}[1][]{
  enhanced, breakable,
  colback=yellow!6, colframe=riskamber,
  boxrule=0.45pt, arc=2pt,
  fonttitle=\bfseries, title={#1},
  left=9pt, right=9pt, top=7pt, bottom=7pt
}
% exhibitbox：表/图外框。left=7pt+right=7pt=14pt 内边距，before upper 内的 \linewidth 已扣除该内边距。
% 约定：\exhibitboxwidth = 真实内边距宽（tcolorbox 已扣好 padding 的 \linewidth）。
% 由于 tabularx/longtable 宽度参数的语义是「外部总宽（含两侧 \tabcolsep）」，内部列实际可用 = W - 2\tabcolsep。
% 所以 tabularx 直接传 \exhibitboxwidth（总宽 == 内边距宽），而其他用途（resizebox、tikz、p{...}）用 \exhibitwidth = \exhibitboxwidth - 2\tabcolsep。
\newlength{\exhibitboxwidth}
\newlength{\exhibitwidth}
\newtcolorbox{exhibitbox}[1][]{
  enhanced, breakable,
  colback=white, colframe=rulegrey,
  borderline north={1.2pt}{0pt}{navy},
  boxrule=0.35pt, arc=1pt,
  fonttitle=\bfseries\color{navy}, title={#1},
  left=7pt, right=7pt, top=6pt, bottom=6pt,
  before upper={\global\exhibitboxwidth=\linewidth\relax
    \global\exhibitwidth=\dimexpr\linewidth-2\tabcolsep\relax},
}
\newtcolorbox{disclosurebox}[1][]{
  enhanced, breakable,
  colback=lightgrey, colframe=steelgrey,
  boxrule=0.35pt, arc=1pt,
  fonttitle=\bfseries, title={#1},
  left=8pt, right=8pt, top=6pt, bottom=6pt
}

% ============ First-page dashboard cover ============
\newcommand{\astockcover}{%
\begin{titlepage}
\thispagestyle{empty}
\begin{tikzpicture}[remember picture, overlay]
  \fill[deepnavy] (current page.north west) rectangle ([yshift=-2.55cm]current page.north east);
  \fill[accentblue] ([yshift=-2.55cm]current page.north west) rectangle ([yshift=-2.66cm]current page.north east);
  \fill[deepnavy] (current page.south west) rectangle ([yshift=1.2cm]current page.south east);
\end{tikzpicture}
\vspace*{0.62cm}
{\color{white}\sffamily\bfseries\small \reportkicker}\\[0.18cm]
{\color{steelgrey}\sffamily\scriptsize \reportscope \quad|\quad \reporttype \quad|\quad \reportdate}\\[1.0cm]
{\fontsize{25}{30}\selectfont\bfseries\color{navy}\reporttitle}\\[0.25cm]
{\fontsize{13}{17}\selectfont\color{steelgrey}\reportsubtitle}\\[0.75cm]
\begin{dashboardbox}[决策摘要]
\begin{tabularx}{\textwidth}{L{3.0cm}X L{3.0cm}X}
\textbf{报告类型} & \reporttype & \textbf{数据截至} & \reportdatacutoff \\
\textbf{覆盖范围} & \reportscope & \textbf{作者} & \reportauthor \\
\end{tabularx}
\vspace{0.45em}
\begin{houseviewbox}[本机构观点]
\reporthouseview
\end{houseviewbox}
\begin{sourcequalitybox}[证据质量]
\reportquality
\end{sourcequalitybox}
\end{dashboardbox}
\vfill
{\scriptsize\color{steelgrey}\reportdisclaimer}
\end{titlepage}
}


\hypersetup{
  pdfauthor={\reportauthor},
  pdftitle={\reporttitle}
}

\newcommand{\qzero}{\textbf{\textcolor{riskred}{q=0 / 仅观察}}}
\newcommand{\blocked}{\textbf{\textcolor{riskred}{未获批准}}}
\newcommand{\watch}{\textbf{\textcolor{riskamber}{观察清单}}}
\newcommand{\sourceid}[1]{\evidencetag{#1}}

\begin{document}

\astockcover

\tableofcontents
\clearpage

\chapter{结论、覆盖范围与当前动作}
\begin{dashboardbox}[一页结论：研究完成的边界，而非配置授权]
\begin{tabularx}{\textwidth}{L{3.3cm}X}
\textbf{覆盖对象} & 3 只 ETF（511010、510880、512890）+ 7 只公募 REIT（508069、508011、508068、508077、508055、508080、508096）；另有“受保障流动性”与“防御分散器”两项未命名类别。\\
\textbf{R0 证据状态} & 已归档产品法律文件、2025 年年报、2026Q2 报告、项目/持仓穿透和上交所 2026-08-12 同日快照；季度财务均明确标为\textbf{未经审计}。\\
\textbf{当前投资动作} & \qzero。\textbf{不新增、不替换、不配置任何 ETF 或 REIT}；并非因为某一项已被否定，而是每一项仍缺执行、当前估值、账户/税务或组合回放所必需字段。\\
\textbf{策略目标口径} & 5\% 是待检验的总回报研究目标，不是年化底线、预测或概率承诺；3\% 仅可指滚动 12 个月、税后费后、实际结算且可复投的净现金，不包括卖出本金、未实现收益、外部入金或借款。\\
\textbf{发布状态} & \textbf{研究框架通过；实施未获批准。} 本备忘录允许继续观察和补证，不允许把历史分派率、标签或单日成交额转化为买入理由。\\
\end{tabularx}
\end{dashboardbox}

\section{本次纠正了什么}

此前的策略框架定义了可能的资产类别与条件性参考权重，但它本身不能等同于底层标的尽调。本案把所有被点名的交易产品逐一放入同一审计框架：\textbf{产品是什么、实际拥有或收费什么、历史现金/收益来自哪里、估值能否成立、主要风险是什么、还缺什么以及能否改变 q。} 结果是十个产品均有可追溯产品卡；两项未命名类别被明确标为“不可研究、不可准入”，不再以类别名称掩盖缺口。

\begin{houseviewbox}[本机构观点]
收益标签不是风险分类。511010 是集中于五年期国债曲线的利率暴露；510880 的已披露 Q2 行业权重集中于采矿与金融；512890 的金融行业权重超过一半；7 只 REIT 则分别承受交通、北京社区商业、租赁住房、亦庄产业园和能源现金流驱动。它们不能仅因历史分派或“低波动/红利”名称被拼接成确定的 5\%/3\% 产品组合。
\end{houseviewbox}

\section{全产品决策矩阵}

\begin{exhibitbox}[表 1：已命名产品逐项处置（截至 2026-08-12）]
\scriptsize
\begin{tabularx}{\exhibitboxwidth}{L{1.35cm}L{2.7cm}L{3.2cm}X L{2.0cm}}
\toprule
\textbf{代码} & \textbf{产品} & \textbf{经济暴露} & \textbf{已完成的核心核验} & \textbf{当前状态} \\
\midrule
511010 & 国泰上证5年期国债ETF & 五年期国债、利率曲线 & 法律条款、费率、年报/Q2 NAV、债券占比及市场快照；缺久期/DV01/同日 IOPV & \qzero \\
510880 & 华泰柏瑞上证红利ETF & 红利股票；采矿、金融权重较高 & 完全复制规则、年报/Q2、分派记录、行业/前十大持仓及快照；缺当前持仓与税后现金 & \qzero \\
512890 & 华泰柏瑞中证红利低波ETF & 红利低波股票；金融权重较高 & 完全复制规则、年报/Q2、持仓/行业与快照；近三年未实施利润分配 & \qzero \\
508069 & 华夏南京交通高速公路REIT & 单项收费公路、车流/收费 & 年报/Q2、通行费、OCF、分派、项目债务线索及快照；缺期限/估值/深度 & \qzero \\
508011 & 嘉实物美消费REIT & 北京四处社区商业、租户租金 & 年报/Q2、出租率/WALT/租户集中、OCF/分派与快照；缺租户信用和债务勾稽 & \qzero \\
508068 & 华夏北京保障房REIT & 北京六处公租房、政策租金 & 年报/Q2、出租/收缴、扩募、OCF/分派与快照；缺扩募后可比估值 & \qzero \\
508077 & 华夏华润有巢租赁住房REIT & 上海三处市场化长租公寓 & 年报/Q2、出租率、租金、短 WALT、扩募与快照；缺续租压力测试 & \qzero \\
508055 & 汇添富上海地产租赁住房REIT & 上海两处保租房 & 年报/Q2、出租率、收缴、分派与快照；首个报告年度不完整 & \qzero \\
508080 & 中金亦庄产业园REIT & 亦庄产业园、工业租户 & 年报/Q2、出租率、最大租户、利润--现金差异与快照；缺现金桥和租约/债务包 & \qzero \\
508096 & 中航京能国际能源REIT & 光伏及水电、发电/电价/应收 & 更名/扩募、年报/Q2、OCF/分派与快照；年报 NAV 多口径未勾稽 & \qzero \\
\bottomrule
\end{tabularx}
\sourcenote{产品宇宙、ETF/REIT 核验包与上交所快照：\sourceid{S04--S28}。所有“缺口”均记录于 source exhaustion log。}
\end{exhibitbox}

\begin{exhibitbox}[图 1：十只已命名产品的经济暴露与当前研究状态]
\centering
\begin{tikzpicture}[
  x=1cm,y=1cm,
  product/.style={rounded corners=2pt, draw=navy, fill=blue!4, minimum width=2.15cm, minimum height=0.62cm, align=center, font=\scriptsize},
  reit/.style={rounded corners=2pt, draw=riskamber!80!black, fill=yellow!10, minimum width=2.15cm, minimum height=0.62cm, align=center, font=\scriptsize},
  exposure/.style={rounded corners=2pt, draw=steelgrey, fill=panelgrey, minimum width=2.35cm, minimum height=0.62cm, align=center, font=\scriptsize},
  blocked/.style={rounded corners=2pt, draw=riskred, fill=red!6, minimum width=2.15cm, minimum height=0.62cm, align=center, font=\scriptsize\color{riskred}},
  arrow/.style={-{Stealth[length=2mm]}, draw=mutedgrey, line width=.35pt}
]
\node[product] (p1) at (0,3.10) {511010\\五年期国债};
\node[product] (p2) at (0,2.20) {510880\\红利权益};
\node[product] (p3) at (0,1.30) {512890\\红利低波权益};
\node[reit] (p4) at (0,0.40) {508069\\收费公路};
\node[reit] (p5) at (0,-0.50) {508011\\社区商业};
\node[reit] (p6) at (0,-1.40) {508068 / 508077 / 508055\\租赁住房风险簇};
\node[reit] (p7) at (0,-2.30) {508080\\产业园租户};
\node[reit] (p8) at (0,-3.20) {508096\\水光能源};
\node[exposure] (e1) at (4.35,3.10) {利率曲线};
\node[exposure] (e2) at (4.35,1.75) {权益因子、金融/周期集中};
\node[exposure] (e3) at (4.35,0.00) {项目运营、租约/车流/资源};
\node[blocked] (b) at (8.30,0.00) {全部 q=0\\不构成配置};
\draw[arrow] (p1.east) -- (e1.west);
\draw[arrow] (p2.east) -- (e2.west);
\draw[arrow] (p3.east) -- (e2.west);
\draw[arrow] (p4.east) -- (e3.west);
\draw[arrow] (p5.east) -- (e3.west);
\draw[arrow] (p6.east) -- (e3.west);
\draw[arrow] (p7.east) -- (e3.west);
\draw[arrow] (p8.east) -- (e3.west);
\draw[arrow] (e1.east) -- (b.north west);
\draw[arrow] (e2.east) -- (b.west);
\draw[arrow] (e3.east) -- (b.south west);
\end{tikzpicture}
\sourcenote{穿透与资料状态：\sourceid{S04--S28}。蓝色为 ETF、黄色为公募 REIT；这是定性经济暴露图，不是相关性、收益预测或配置建议。}
\end{exhibitbox}

\section{未命名类别不是已尽调产品}

\begin{riskbox}[两个不可跳过的范围缺口]
\textbf{受保障流动性类别：}当前 CNY200,000 的防御性停泊表述没有指定存款、直接国债或货币市场产品，因而不能核验法律债权、存款保险合并口径、期限、提前支取、税费或账户资格。它只是停泊\emph{类别}，不可以被称为已选产品或已知收益率。\par
\textbf{防御分散器类别：}没有标的、非杠杆筛选规则、相关性、压力测试和流动性资料，不能假设其能够降低组合回撤，更不能为 3\% 现金目标贡献一个虚构收益率。两类均维持 q=0，先冻结有限候选清单，才有下一轮产品尽调。
\end{riskbox}

\section{从观察到可讨论配置：必须同时满足的八个门槛}

\begin{exhibitbox}[图 2：从产品研究到可讨论配置的八项串联门槛]
\centering
\begin{tikzpicture}[
  gate/.style={rounded corners=2pt, draw=navy, fill=panelgrey, minimum width=2.18cm, minimum height=0.82cm, align=center, font=\scriptsize},
  finalgate/.style={rounded corners=2pt, draw=riskred, fill=red!6, minimum width=2.18cm, minimum height=0.82cm, align=center, font=\scriptsize\color{riskred}},
  arrow/.style={-{Stealth[length=2mm]}, draw=mutedgrey, line width=.4pt}
]
\node[gate] (g1) at (0,1.20) {1\\法律身份与\\资产清单};
\node[gate] (g2) at (2.75,1.20) {2\\同日市场价与\\NAV / IOPV};
\node[gate] (g3) at (5.50,1.20) {3\\点差、深度与\\20 日容量};
\node[gate] (g4) at (8.25,1.20) {4\\账户、税务、\\结算与适配};
\node[gate] (g5) at (1.38,-0.40) {5\\久期/持仓或\\项目运营与债务};
\node[gate] (g6) at (4.13,-0.40) {6\\共同风险、\\地域与管理人压力};
\node[gate] (g7) at (6.88,-0.40) {7\\历史 TWRR 与\\已结算净现金回放};
\node[finalgate] (g8) at (9.63,-0.40) {8\\独立风险预算与\\组合层批准};
\draw[arrow] (g1.east) -- (g2.west);
\draw[arrow] (g2.east) -- (g3.west);
\draw[arrow] (g3.east) -- (g4.west);
\draw[arrow] (g4.south) |- (g5.east);
\draw[arrow] (g5.east) -- (g6.west);
\draw[arrow] (g6.east) -- (g7.west);
\draw[arrow] (g7.east) -- (g8.west);
\end{tikzpicture}
\sourcenote{门禁来自本案产品宇宙、回报模型、风险框架和资料穷尽日志。八项必须同时满足；图中顺序仅表达依赖关系，不表示任何产品将获准配置。}
\end{exhibitbox}

\begin{exhibitbox}[图 3：历史资料、市场快照与当前实施资料之间的断点]
\centering
\begin{tikzpicture}[
  x=1cm,y=1cm,
  observed/.style={rounded corners=2pt, draw=navy, fill=blue!4, minimum width=2.55cm, minimum height=1.02cm, align=center, font=\scriptsize},
  missing/.style={rounded corners=2pt, draw=riskred, fill=red!6, minimum width=2.70cm, minimum height=1.02cm, align=center, font=\scriptsize\color{riskred}},
  arrow/.style={-{Stealth[length=2mm]}, draw=mutedgrey, line width=.4pt},
  note/.style={font=\scriptsize\color{steelgrey}, align=center}
]
\node[observed] (annual) at (-5.05,0) {2025-12-31\\经审计年报\\历史 NAV / 现金};
\node[observed] (q2) at (-1.68,0) {2026-06-30\\未经审计 Q2\\持仓 / 项目运营};
\node[observed] (market) at (1.68,0) {2026-08-12 16:29\\上交所快照\\价格 / 成交额};
\node[missing] (current) at (5.15,0) {当前仍缺\\同日 NAV / IOPV、点差\\深度、容量与账户资料};
\draw[arrow] (annual.east) -- (q2.west);
\draw[arrow] (q2.east) -- (market.west);
\draw[arrow] (market.east) -- (current.west);
\node[note] at (0,-1.15) {日期错配与缺失字段意味着：不能将历史 NAV 诊断升级为当前估值，\\也不能将单日成交额升级为可执行容量。};
\end{tikzpicture}
\sourcenote{年报/Q2 与交易所快照：\sourceid{S04--S27}。所有十项的文件级日期、相对路径与 SHA-256 见 data/legal\_document\_manifest\_20260812.json；红色节点为实施阻断项。}
\end{exhibitbox}

未满足任何一项时，研究降级为观察；不以高收益替代、杠杆、降低信用质量或出售本金来“修复”目标。


\chapter{方法、证据层级与估值边界}
\section{研究问题与边界}

本报告回答的是一个有限问题：策略中\textbf{所有已命名产品}是否已经被逐项调查到足以描述经济实质和识别实施缺口。它不回答“今天买哪一只”，也不假定旧框架中的条件性 10\%/15\%/40\%/25\%/10\% 参考向量是有效配置。活动模型仍是未命名保护性流动性类别 100\%，全部市场产品权重为零。

本案采用自定义的“收益型 ETF / 公募 REIT / 流动性工具”覆盖包，而不是产业链报告模板。每一项产品必须经过：法律与管理结构、底层穿透、经审计年度事实、未经审计季度更新、历史现金/回报、价格与估值可比性、流动性/账户可执行性、共同风险、升级触发条件八个模块。直接股票、杠杆、衍生品、低信用收益替代品均不在本案宇宙内。

\section{证据层级与时间标签}

\begin{exhibitbox}[表 2：本案证据层级]
\small
\begin{tabularx}{\exhibitboxwidth}{L{2.0cm}L{3.1cm}X}
\toprule
\textbf{层级} & \textbf{来源} & \textbf{可以支持 / 不可以支持} \\
\midrule
一级、经审计年度 & 2025 年报、审计意见、法律文件 & 历史财务、NAV、披露的资产/持仓和年度分派；不能代表 2026-08-12 的资产价值或未来现金。\\
一级、未经审计季度 & 2026Q2 定期报告、运营公告 & 阶段性运营、持仓和管理人已披露分派口径；不能升级为经审计事实，季度分派不能机械年化。\\
一级、交易所快照 & 上交所 2026-08-12 16:29:05 快照 & 同日最后价、成交量、成交额、OHLC；不能证明点差、深度、容量、做市质量、同日 IOPV/NAV、税费或可执行性。\\
二/三级、卖方资料 & 原始 PDF、转载、预览、检索失败日志 & 仅做语料完整性审计。本案没有可用的产品级评级、目标、预测和方法论锚；所有卖方权重为零。\\
\bottomrule
\end{tabularx}
\sourcenote{一级来源索引：\sourceid{S01--S28}；卖方穷尽记录见 broker/street consensus packet。}
\end{exhibitbox}

时间错配会直接改变结论：ETF 的净值和持仓多为 2026-06-30，市场快照为 2026-08-12；REIT 的可用每份净资产为 2025-12-31 或不同披露口径。故本文将“价格相对历史 NAV”称为\textbf{日期错配诊断}，不用“当前折溢价”，不发布目标价或公允价值区间。

\section{现金与总回报：不双重计算的定义}

\begin{houseviewbox}[唯一允许的未来验证口径]
\textbf{净现金收益率} = 在滚动 12 个月内实际结算的合格利息、ETF 分派或 REIT 现金分派，扣除税费、领取费用、持续执行成本和已记录的分派下调后，除以期初组合资本。卖出资产所得、未实现市值变动、NAV 增长、外部入金、借款及未验证的返还本金均不得计入。\par
\textbf{总回报}采用外部现金流分段的 TWRR；实际投资者另行报告 MWRR/XIRR。合格现金只在总回报中计入一次，复投只改变下一期资本基数。没有逐点历史产品集合、成本税费、容量假设和基准对照的回放之前，5\% 和 3\% 均不可被称为已验证结果。
\end{houseviewbox}

\section{估值和执行为何必须拆开}

产品在交易所成交并不代表交易价格就是可证明的价值，交易额也不等于可部署容量。本案把以下三项分别记录：
\begin{enumerate}
  \item \textbf{市场观察：}2026-08-12 的最后价和成交额；
  \item \textbf{历史 NAV 诊断：}仅当历史每份 NAV 有单一可引用口径时，计算 $P/NAV-1$，同时展示日期差；
  \item \textbf{执行资格：}同日 NAV/IOPV、点差、深度、容量、账户/税务/结算和风险预算均完成后才可判断。
\end{enumerate}
508096 的年报出现三个未勾稽的每份 NAV 相关数字，连第 2 项诊断也被阻断。对其余九只，诊断不构成“贵/便宜”结论，原因是资产、债务、租约、利率或持仓可能已经变化。

\section{配置前的产品层与组合层门槛}

产品层通过不等于组合层通过。即使某只 REIT 补齐估值，REIT 整体仍至少需要五个独立合格项目、三个以上管理人/赞助人组和三个以上运营现金流行业；单项目不超过总组合 5\%，全部 REIT 不超过 25\%。红利 ETF 之间的金融与周期风险、三只租赁住房 REIT 的同簇风险、地域和管理人关联均需要共同风险限制。当前没有批准的风险预算，因此所有非零权重仍被锁定。


\chapter{利率暴露产品：511010 国债 ETF}
\section{产品定义：利率暴露，不是保护性现金}

511010（国泰上证5年期国债ETF）由国泰基金管理有限公司管理、中国建设银行托管，以优化抽样跟踪上证5年期国债指数收益率。产品资料概要披露，建仓完成后不低于 90\% 的基金资产应投资于标的指数成份券或备选成份券；管理费 0.15\%，托管费 0.05\%，年化运营费用估算 0.21\%。因此其经济本质是\textbf{集中于五年期国债曲线的 ETF 利率敞口}，不是到期本金确定的存款，也不是没有净值波动的“保护性流动性”。

\begin{exhibitbox}[表 3：511010 的已核验事实与未完成模块]
\small
\begin{tabularx}{\exhibitboxwidth}{L{3.4cm}X}
\toprule
\textbf{模块} & \textbf{已核验事实 / 结论} \\
\midrule
2025 年报（经审计） & 基金资产净值 CNY3.021bn；每份 NAV CNY140.027；年度 NAV 回报 0.59\%；现金分配 CNY21.042/10 份，合计 CNY38.003m。历史分配不等于未来收益。\\
2026Q2（未经审计） & 2026-06-30 NAV CNY5.117bn、每份 NAV CNY140.753、三个月 NAV 回报 0.78\%；国债占 NAV 98.89\%。\\
穿透集中度 & Q2 前五项国债占 NAV 87.30\%，最大单项“25附息国债20”占 20.75\%。这是成份券集中度描述，不是久期风险模型。\\
交易所观察 & 2026-08-12 16:29:05 最后价 CNY141.084，当日成交额 CNY1.345bn。快照没有同日 IOPV/NAV、点差、深度或容量。\\
估值状态 & 价格相对 2026-06-30 NAV 的日期错配诊断约 +0.24\%；相隔 43 天，不能称为当前溢价或公允价值。\\
行动 & \qzero；任何条件性 15\% 参考权重都未获批准。\\
\bottomrule
\end{tabularx}
\sourcenote{年报、Q2、产品资料概要、持仓表和上交所快照：\sourceid{S04--S06, S27, S28}。}
\end{exhibitbox}

\section{真正的风险变量}

该产品的风险问题不是“国债是否安全”这一抽象判断，而是：当收益率曲线平移或陡峭化时，基金的修正久期、DV01 和凸性是多少；优化抽样相对指数如何偏离；当前成交价格是否可以以给定规模执行；以及账户、税费与申赎机制是否合规。现有源包没有这些当前字段。即使披露资产几乎全为国债，也不能把它改写为无风险现金或替代未命名的保护性流动性类别。

\begin{riskbox}[升级触发条件]
在 q 能从零变动之前，需要同一时点的基金 NAV/IOPV、完整成份券与修正久期/DV01/凸性、至少 20 个交易日的点差/深度/冲击成本、申赎状态以及账户和税费说明。组合层还需测试利率冲击与其他风险资产的共同回撤。没有这些资料，历史 0.59\% 年度 NAV 回报和 0.78\% 季度回报都不能外推。
\end{riskbox}

\section{对 3\% 现金目标的含义}

2025 年度现金分配只能证明该年度存在一次已披露分配，不能证明持有人未来可获得相同金额，更不能直接计入一个 3\% 的可复投净现金目标。只有实际结算、扣除税费和成本且法律性质已确认的现金，才可进入后续的净现金台账；其余市值或 NAV 变化只属于总回报的一部分。


\chapter{权益现金流产品：510880 与 512890}
\section{510880：红利标签背后的周期与金融权重}

510880（华泰柏瑞上证红利ETF）完全复制上证红利指数，产品资料概要规定一般不低于 95\% 的基金资产投向指数成份股和备选成份股；管理费 0.50\%、托管费 0.10\%、年化运营费用估算 0.60\%。它是权益 ETF，而不是现金流合同。2026Q2 的证券资产占总资产 98.73\%，已披露行业中采矿业占 NAV 33.95\%、金融业占 31.34\%，前十持仓占 30.95\%。这意味着“红利”并没有消除周期、商品和金融因子集中度。

\begin{exhibitbox}[表 4：510880 证据卡]
\small
\begin{tabularx}{\exhibitboxwidth}{L{3.35cm}X}
\toprule
\textbf{项目} & \textbf{核验结果} \\
\midrule
2025 年报（经审计） & NAV CNY19.265bn、每份 NAV CNY3.1813、年度 NAV 回报 -0.45\%。2025/2024/2023 每 10 份现金分配分别为 CNY1.4200/CNY1.3100/CNY1.3800。\\
2026Q2（未经审计） & NAV CNY22.276bn、每份 NAV CNY2.9103、三个月 NAV 回报 -10.36\%。行业与持仓表仅代表 2026-06-30。\\
代表性持仓 & 中远海控、中谷物流、潞安环能、兖矿能源、广汇能源、山煤国际、淮北矿业、平煤股份、陕西煤业和西部矿业。该列表不是当前组合。\\
交易所观察 & 2026-08-12 最后价 CNY3.303，成交额 CNY389.653m；不证明可按计划规模成交。\\
历史 NAV 诊断 & 最后价相对 2026-06-30 每份 NAV 约 +13.49\%，相隔 43 天；仅是日期错配诊断，不能称为当前溢价、目标价或估值结论。\\
行动 & \qzero；40\% 条件性参考权重没有资格转为配置。\\
\bottomrule
\end{tabularx}
\sourcenote{年报、Q2、产品资料概要、指数方法与上交所快照：\sourceid{S07--S09, S27, S28}。}
\end{exhibitbox}

\section{512890：低波动名称不等于低组合风险或现金}

512890（华泰柏瑞中证红利低波ETF）完全复制中证红利低波动指数，资料概要规定一般不低于 90\% NAV、且非现金资产中不低于 80\% 投于成份股/备选成份股；管理费 0.50\%、托管费 0.10\%、年化运营费用估算 0.60\%。2026Q2 的证券资产占总资产 97.77\%，金融业占 NAV 54.13\%，前十持仓占 25.01\%。它以银行为主的因子敞口很清晰，不可因“低波”三个字假定其与 510880 或租金类 REIT 独立。

\begin{exhibitbox}[表 5：512890 证据卡]
\small
\begin{tabularx}{\exhibitboxwidth}{L{3.35cm}X}
\toprule
\textbf{项目} & \textbf{核验结果} \\
\midrule
2025 年报（经审计） & NAV CNY26.754bn、每份 NAV CNY1.1745、年度 NAV 回报 4.74\%。年报表述过去三年没有实施利润分配。\\
2026Q2（未经审计） & NAV CNY31.290bn、每份 NAV CNY1.0668、三个月 NAV 回报 -10.83\%。金融业权重 54.13\%。\\
代表性持仓 & 上海银行、南京银行、平安银行、成都银行、格力电器、江苏银行、山东高速、沪农商行、陕鼓动力和民生银行。\\
交易所观察 & 2026-08-12 最后价 CNY1.159，成交额 CNY517.864m；同样缺少 IOPV、点差和深度。\\
历史 NAV 诊断 & 最后价相对 2026-06-30 每份 NAV 约 +8.64\%，相隔 43 天；不是当前折溢价。\\
行动 & \qzero；不得把它当作未达 3\% 现金目标时的自动替代品。\\
\bottomrule
\end{tabularx}
\sourcenote{年报、Q2、产品资料概要与上交所快照：\sourceid{S10--S11, S27, S28}。}
\end{exhibitbox}

\section{两只 ETF 的共同门禁}

两只权益 ETF 的历史现金分配与基金层利润分配政策不同，不能混为“股息率”。下一轮需要同日全持仓文件、发行人集团重叠、TTM 实际结算现金、税后/费用处理、指数调整规则、跟踪差、压力期间回撤及相对 510880/512890/REIT 的共同风险。没有这些，任何“红利现金贡献”的数字都只是未经验证的假设。


\chapter{公募 REIT：收费公路与社区商业}
\section{508069：单资产收费公路，关键在车流、期限与融资}

508069（华夏南京交通高速公路REIT）穿透至南京绕越高速公路东南段一项资产，现金流以通行费为主并有少量租赁收入。2025 年报（经审计）披露收入 CNY490.114m、净利润 CNY94.493m、经营现金流 CNY331.088m、分派率 9.19\%、每份 NAV CNY4.8970。2026Q2（未经审计）披露收入 CNY113.090m、经营现金流 CNY61.292m，项目公司通行费收入同比 -8.20\%；管理人列示的年化率 9.75\% 仅是发行人口径的历史算术，不是本报告预测。

\begin{exhibitbox}[表 6：508069 产品卡]
\small
\begin{tabularx}{\exhibitboxwidth}{L{3.35cm}X}
\toprule
\textbf{维度} & \textbf{核验与风险含义} \\
\midrule
现金流和运营 & Q2 可分配金额 CNY59.460m、实际分派 CNY164.400m；收入、费用和收现并非线性发生，不能将季度口径机械年化。\\
资产与融资线索 & 项目公司披露长期应付款 CNY1.817bn、长期银行借款 CNY263.7m及一年内到期部分 CNY16.5m。未完成基金--ABS--项目公司的法律/会计勾稽，故不计算“基金杠杆率”。\\
市场观察 & 2026-08-12 最后价 CNY5.221、成交额 CNY0.347m。价格相对 2025-12-31 NAV 的历史诊断约 +6.62\%，日期相隔 224 天，绝非当前 P/NAV。\\
不可省略的尽调 & 车流、收费标准、路网分流、剩余收费权、资本开支、长期应付款利率/期限/偿还顺位和再融资情景。\\
行动 & \qzero；单项目也受总组合 5\% 项目上限约束，且只有在组合层门禁通过后才可讨论。\\
\bottomrule
\end{tabularx}
\sourcenote{508069 年报、Q2、运营披露与上交所快照：\sourceid{S12--S13, S27, S28}。}
\end{exhibitbox}

\section{508011：北京社区商业，关键在租户、资产分化与结构勾稽}

508011（嘉实物美消费REIT）拥有北京丰台、海淀和西城的四处社区商业资产，现金流来自商业租赁与物业管理。2025 年报（经审计）披露收入 CNY105.051m、净利润 CNY31.350m、经营现金流 CNY77.270m、分派率 4.38\%、每份 NAV CNY2.3089。2026Q2（未经审计）组合出租率 93.66\%、加权剩余租期 5.94 年，前五租户占租金收入 53.02\%；大成和华天出租率分别为 89.32\% 和 90.93\%，呈现项目间差异。

\begin{exhibitbox}[表 7：508011 产品卡]
\small
\begin{tabularx}{\exhibitboxwidth}{L{3.35cm}X}
\toprule
\textbf{维度} & \textbf{核验与风险含义} \\
\midrule
现金流和运营 & Q2 收入 CNY27.424m、经营现金流 CNY12.980m、可分配金额 CNY18.114m、实际分派 CNY35.240m；管理人列示年化率 5.03\% 不是前瞻收益率。\\
结构线索 & 三个项目公司分别披露长期借款 CNY438.612m、CNY107.881m、CNY88.490m；聚合项目公司资产 CNY700.154m、负债 CNY752.178m。该汇总口径不能直接推成基金层杠杆。\\
市场观察 & 2026-08-12 最后价 CNY3.460、成交额 CNY0.876m。相对 2025 年末 NAV 的历史诊断约 +49.86\%，但资产运营、估值和资本结构已有时点差。\\
不可省略的尽调 & 匿名大租户信用、租约/续约、低出租率资产的有效租金和资本开支、项目公司负债与基金净资产的结构桥。\\
行动 & \qzero；唯一归档的原始券商 PDF也不含现行评级、目标、预测或估值法，卖方权重仍为零。\\
\bottomrule
\end{tabularx}
\sourcenote{508011 年报、Q2、运营披露、上交所快照和卖方语料审计：\sourceid{S14--S15, S27, S28}。}
\end{exhibitbox}

\begin{riskbox}[REIT 分派的统一警告]
公募 REIT 合同中的“符合条件时年度分配不低于可供分配金额 90\%”是分派政策，不是固定收益率或总回报下限。所有七份 Q2 报告均未审计，且发行人提示收入和费用并不均匀，季度可分配金额、实际分派和管理人年化率不能被线性外推。
\end{riskbox}


\chapter{公募 REIT：北京与上海租赁住房}
\section{租赁住房不是单一“防御资产”：三只产品共享租金、续约与政策簇}

508068、508077 与 508055 均以租赁住房为核心，但收入机制、地域、价格约束、扩募历史与租约风险并不相同。三者不能因为同属“住房 REIT”而被假定为独立，亦不能把较高的期末出租率等同于可持续分派。三只 Q2 数据均为\textbf{未经审计}，且扩募或不足完整年度会破坏同比和年化比较。

\section{508068：北京保障房，稳定性与提价弹性受限并存}

508068（华夏北京保障房REIT）穿透至北京六处公共租赁住房，现金流为政府指导租金。2025 年报（经审计）披露收入 CNY104.847m、净利润 CNY45.376m、经营现金流 CNY72.049m、分派率 2.53\%、每份 NAV CNY2.9927。Q2 披露合计 4,264 套、可出租面积 241,509.97 平方米、出租率 94.07\%、收缴率 96.30\%，且披露租金标准约为 CNY35--65/平方米/月、一个租期内不调整。

\begin{exhibitbox}[表 8：508068 产品卡]
\small
\begin{tabularx}{\exhibitboxwidth}{L{3.35cm}X}
\toprule
\textbf{维度} & \textbf{核验与风险含义} \\
\midrule
Q2 经营/分派 & 收入 CNY32.292m、经营现金流 CNY24.876m、可分配金额 CNY22.430m、实际分派 CNY43.453m；管理人列示年化率 3.13\% 仅为历史披露。\\
扩募与结构 & 2025 年扩募增加资产，资产基数发生变化；两个项目公司聚合资产 CNY1.973bn、负债 CNY1.690bn，燕保宜居长期应付款 CNY836.340m。不能由此直接计算基金层杠杆。\\
市场/NAV 诊断 & 2026-08-12 最后价 CNY3.863、成交额 CNY3.198m；相对 2025-12-31 NAV 的历史诊断约 +29.08\%，相隔 224 天，非当前估值。\\
关键补证 & 政策指导租金/准入/补贴变化、扩募前后同口径现金流、项目级空置与收缴、债务期限和资本开支。\\
行动 & \qzero；不因“保障房”名称把政策稳定性误写为固定现金流。\\
\bottomrule
\end{tabularx}
\sourcenote{508068 年报、Q2、法律文件与上交所快照：\sourceid{S16--S17, S27, S28}。}
\end{exhibitbox}

\section{508077：市场化上海长租，短租期与租金下行是核心变量}

508077（华夏华润有巢租赁住房REIT）持有上海泗泾、东部经开区及马桥三处市场化长租公寓，收入还包括停车和增值服务。2025 年报（经审计）披露收入 CNY79.052m、净利润 CNY11.480m、经营现金流 CNY52.832m、分派率 1.80\%、每份 NAV CNY2.3431。Q2 披露按面积出租率 94.98\%、平均租金 CNY2.12/平方米/天，同比 -2.30\%，加权剩余租期仅 0.49 年；高收缴率 99.92\% 不能取代续租和再出租压力测试。

\begin{exhibitbox}[表 9：508077 产品卡]
\small
\begin{tabularx}{\exhibitboxwidth}{L{3.35cm}X}
\toprule
\textbf{维度} & \textbf{核验与风险含义} \\
\midrule
Q2 经营/分派 & 收入 CNY34.868m、经营现金流 CNY21.148m、可分配金额 CNY22.767m、实际分派 CNY22.539m；管理人列示年化率 3.48\% 不是预测。\\
扩募与结构 & 2025 年末扩募使可比资产基数改变；项目公司聚合资产 CNY1.630bn、负债 CNY1.700bn，两个资产披露长期应付款。必须先做结构勾稽。\\
市场/NAV 诊断 & 2026-08-12 最后价 CNY2.666、成交额 CNY15.356m；相对 2025 年末 NAV 的历史诊断约 +13.78\%，不等于当前价格结论。\\
关键补证 & 续租率、有效租金、免租优惠、各项目新签/续签节奏、扩募后同口径经营序列、债务和资本开支。\\
行动 & \qzero；租金同比下降和短 WALT 不能只用期末出租率覆盖。\\
\bottomrule
\end{tabularx}
\sourcenote{508077 年报、Q2、法律文件与上交所快照：\sourceid{S18--S19, S27, S28}。}
\end{exhibitbox}

\section{508055：上海保租房，短历史与资产内部分化要先处理}

508055（汇添富上海地产租赁住房REIT）包括江月路和虹桥两处保障性租赁住房，现金流来自租金、配套商业和停车。2025 年报是从 3 月 21 日起的部分首年，披露收入 CNY63.123m、净利润 CNY16.626m、经营现金流 CNY51.074m、分派率 2.90\%、每份 NAV CNY2.6811；这一期间不能视为完整年度基准。Q2 组合出租率 96.43\%、平均租金 CNY2.79/平方米/天、收缴率 99.32\%，但江月路出租率 97.69\% 高于虹桥的 93.16\%，需要资产级而非组合平均数判断。

\begin{exhibitbox}[表 10：508055 产品卡]
\small
\begin{tabularx}{\exhibitboxwidth}{L{3.35cm}X}
\toprule
\textbf{维度} & \textbf{核验与风险含义} \\
\midrule
Q2 经营/分派 & 收入 CNY18.248m、经营现金流 CNY13.311m、可分配金额 CNY13.899m、实际分派 CNY25.770m；管理人列示年化率 3.27\% 是历史口径。\\
结构线索 & 项目公司聚合资产 CNY1.302bn、负债 CNY1.020bn；江月路与虹桥长期应付款分别为 CNY636.000m 与 CNY257.223m。未做基金层杠杆判断。\\
市场/NAV 诊断 & 2026-08-12 最后价 CNY3.304、成交额 CNY0.829m；相对 2025 年末 NAV 的历史诊断约 +23.23\%，但历史较短且时点错配。\\
关键补证 & 完整年度经营序列、虹桥资产的租金/出租率与资本开支、长期应付款条款、项目级税后现金。\\
行动 & \qzero；不以部分首年年化或一次分派推导长期可复投现金。\\
\bottomrule
\end{tabularx}
\sourcenote{508055 年报、Q2、法律文件与上交所快照：\sourceid{S20--S21, S27, S28}。}
\end{exhibitbox}

\section{租赁住房簇的组合含义}

三只产品共同暴露于出租率、有效租金、续约、收缴、区域供给、政策和长期应付款等变量。它们必须被视为一个共同风险簇，不能在没有相关性、压力测试和管理人/赞助人关系审查的情况下分散计数。即便单只产品完成补证，REIT 总体仍要满足五个独立项目、三个管理人/赞助人组和三个运营现金流行业的激活条件。


\chapter{公募 REIT：产业园与国际能源}
\section{508080：产业园的租户集中、去化和利润--现金桥}

508080（中金亦庄产业园REIT）穿透至北京经济技术开发区两处面向高端汽车/新能源零部件的产业园，收入来自厂房租赁、物业管理与综合服务。2025 年报为自 2025-06-16 起的部分首年，披露收入 CNY61.901m、净利润 CNY5.002m、经营现金流 CNY57.522m、分派率 2.43\%、每份 NAV CNY2.7325。Q2 显示出租率 87.74\%、平均租金 CNY2.62/平方米/天、收缴率 90.11\%，最大租户占租金 44.04\%。这不是一个可忽略的单一客户风险。

\begin{exhibitbox}[表 11：508080 产品卡]
\small
\begin{tabularx}{\exhibitboxwidth}{L{3.35cm}X}
\toprule
\textbf{维度} & \textbf{核验与风险含义} \\
\midrule
Q2 经营/分派 & 收入 CNY24.166m、净利润 -CNY1.761m、经营现金流 CNY19.030m、可分配金额 CNY17.111m、实际分派 CNY36.000m；管理人列示年化率 5.90\% 不能充当预测。\\
应做的现金桥 & 负利润、正经营现金流和可分配金额并非自动矛盾，但需要租金确认、折旧/公允价值、营运资本、融资、资本开支及实际分派的逐项桥接。现有包未完成。\\
结构线索 & 项目公司资产 CNY1.035bn、负债 CNY736.669m，投资性房地产 CNY935.773m、长期应付款 CNY649.091m；不可直接换算基金层杠杆。\\
市场/NAV 诊断 & 2026-08-12 最后价 CNY2.405、成交额 CNY4.456m；相对 2025 年末 NAV 的历史诊断约 -11.99\%。时间错配和资产变化使其不构成“低估”判断。\\
行动 & \qzero；先验证最大租户、租约、去化、现金桥和债务到期，再讨论任何价值判断。\\
\bottomrule
\end{tabularx}
\sourcenote{508080 年报、Q2、法律文件与上交所快照：\sourceid{S22--S23, S27, S28}。}
\end{exhibitbox}

\section{508096：当前是国际能源 REIT，不再是“光伏单资产”}

508096 的当前名称为\textbf{中航京能国际能源REIT}。上交所 2026-02-13 公告确认更名生效；底层不再是光伏单一叙事，而是陕西榆林、湖北随州两项光伏及云南腾冲两项水电资产。2025 年扩募、2026 年初新增份额后，资产基础已经变化。现金流取决于发电资源、电价、结算、输电/限电、补贴应收和保理，而不是传统商业/住房租金。

\begin{exhibitbox}[表 12：508096 产品卡]
\small
\begin{tabularx}{\exhibitboxwidth}{L{3.35cm}X}
\toprule
\textbf{维度} & \textbf{核验与风险含义} \\
\midrule
2025 年报（经审计） & 收入 CNY307.501m、净利润 CNY100.890m、经营现金流 CNY239.235m、分派率 5.33\%。扩募与资产变化降低与未来期间的直接可比性。\\
Q2（未经审计） & 收入 CNY167.554m、净利润 CNY55.655m、经营现金流 CNY93.771m、可分配金额 CNY125.623m、实际分派 CNY210.074m；管理人列示年化率 10.12\% 是历史口径，受经营波动和保理影响。\\
NAV 阻断项 & 年报同时出现 CNY9.1720、CNY8.1264、CNY8.8388 三个每份 NAV 相关数字，口径/时点尚未勾稽。故不计算价格/NAV，即使 2026-08-12 最后价为 CNY9.264。\\
融资/运营线索 & 项目公司聚合资产 CNY4.747bn、负债 CNY3.470bn；湖北晶泰和江山永宸披露短借，后者同比增加 36.07\%。需单独评估补贴应收、保理、借款及来水/光照/限电敏感性。\\
市场观察 & 2026-08-12 成交额 CNY3.043m；没有同日 NAV、点差、深度或容量，因此不产生可执行性结论。\\
行动 & \qzero；NAV 勾稽是估值讨论的立即阻断项。\\
\bottomrule
\end{tabularx}
\sourcenote{508096 年报、Q2、运营/更名披露、法律文件与上交所快照：\sourceid{S24--S26, S27, S28}。}
\end{exhibitbox}

\section{两者的共同误区}

产业园和新能源的“高历史年化分派口径”不能承担策略现金目标。前者需要租户和会计--现金桥，后者需要资源、电价、应收与保理桥；两者都需要同时点估值和实施流动性。不同的现金流来源提升了研究复杂度，不自动提供组合分散化。


\chapter{组合层现金、重叠风险与实施门禁}
\section{现金目标：当前不能证明 3\% 净现金}

本案没有把任何历史产品分派当作未来现金预测。三个 ETF 的分配政策和历史记录不同；七只 REIT 都有“在符合条件时年度分配不低于可供分配金额 90\%”的政策，但这不是固定收益率。特别是，季度收入、费用、项目收现、扩募、保理和实际分派并不均匀，发行人已提示季度可分配金额不能简单年化。

\begin{exhibitbox}[表 13：历史现金观察不能替代 3\% 目标验证]
\scriptsize
\begin{tabularx}{\exhibitboxwidth}{L{1.25cm}L{2.75cm}L{3.25cm}X}
\toprule
\textbf{代码} & \textbf{可见历史现金资料} & \textbf{为何不能外推} & \textbf{进入净现金模型前仍需补齐} \\
\midrule
511010 & 2025 年现金分配 CNY21.042/10 份 & 单年度历史分配，不是票息承诺 & 实际结算、税费、申赎与未来分派资料 \\
510880 & 2023--2025 均有每 10 份分配记录 & 指数成份、利润、基金决定与税费会变化 & TTM 已结算现金、完整持仓、税后现金桥 \\
512890 & 年报称近三年未实施利润分配 & “红利低波”不代表基金持有人现金 & 基金分配政策、TTM 已结算现金与持仓质量 \\
508069--508096 & 年报历史分派率与 Q2 管理人年化口径 & Q2 未审计；季度现金、扩募、费用、保理和项目运营非线性 & 每项的可分配金额桥、实际结算、税费、债务/资本开支和压力情景 \\
\bottomrule
\end{tabularx}
\sourcenote{ETF 与 REIT 年报/Q2 核验包：\sourceid{S04--S25}。}
\end{exhibitbox}

合格的现金口径必须为滚动 12 个月的税后、费后、已结算、可复投现金；出售资产、未实现市值变动、NAV 增长、外部资金、借款以及未验证的返还本金均排除。若短缺，策略必须记录短缺，而不是以杠杆、低信用替代、被迫购入 REIT 或卖出本金修补。当前没有足够数据将任何产品的历史分派加总为组合 3\%。

\section{5\% 总回报目标：没有可复现实证就不是策略结果}

{\sloppy
总回报应以逐点产品集合和外部现金流分段的 TWRR 计算，现金只计入一次；投资者体验另以 MWRR 与 XIRR 报告。当前缺少历史成份/项目状态、分派、公司行为、再平衡、交易成本、税费、流动性/容量和基准序列，因而没有回放，也没有“每年 5\%”的证据。正式测试至少要与保护性流动性、期限匹配国债暴露和相关多资产替代方案比较，并报告回撤、命中率、换手、容量和现金短缺频率。\par}

\section{共同风险矩阵：名称不同，不等于独立}

\begin{exhibitbox}[表 14：跨产品共同风险与当前控制]
\small
\begin{tabularx}{\exhibitboxwidth}{L{2.3cm}L{4.0cm}X}
\toprule
\textbf{风险簇} & \textbf{成员与传导机制} & \textbf{当前控制/缺口} \\
\midrule
利率与估值 & 511010 的国债久期；所有 REIT 的融资成本和估值折现率 & 缺久期/DV01、同日 NAV/估值桥、债务到期表和联动压力测试。\\
红利/金融/周期权益 & 510880 的采矿/金融权重与 512890 的金融权重 & 需完整同日持仓、发行人集团重叠、因子回撤与现金质量分析。\\
租赁住房 & 508068、508077、508055：出租率、租金、续约、收缴、区域供给、政策/长期应付款 & 三只不得当作三份独立现金来源；需共同风险上限和项目级压力测试。\\
单资产或集中租户 & 508069 单公路、508080 最大租户 44.04\%、508011 前五租户 53.02\% & 需运营 KPI、租约/信用、续约、资本开支及债务资料。\\
能源/资源现金流 & 508096 的水光资源、电价、调度、补贴应收与保理 & 需发电、电价、应收、保理、借款和 NAV 口径勾稽。\\
市场可执行性 & 全部十项的点差、深度、容量、做市、账户和税务 & 单日成交额不等于可执行容量；所有 q 维持为零。\\
\bottomrule
\end{tabularx}
\sourcenote{ETF/REIT 穿透图谱和上交所快照：\sourceid{S04--S28}。}
\end{exhibitbox}

\section{由证据备忘录到实施的门禁表}

\begin{dashboardbox}[实施资格：全部阻断]
\begin{tabularx}{\textwidth}{L{3.3cm}X L{2.7cm}}
\textbf{回报证据} & 未完成固定产品集合的逐点 TWRR/已结算净现金回放与基准比较。 & \blocked \\
\textbf{风险预算} & 未批准共同因子、压力损失、集中度、流动性容量及回撤行动规则。 & \blocked \\
\textbf{适配性与账户} & 未提供账户、税务、结算、托管、投资者适配和资金用途资料。 & \blocked \\
\textbf{执行} & 缺同日 NAV/IOPV、点差、盘口深度、20 日容量、冲击成本和市场做市资料。 & \blocked \\
\textbf{产品数量 q} & 所有 ETF、REIT 与未命名类别均为零；不存在实时执行授权。 & \qzero \\
\end{tabularx}
\end{dashboardbox}

\section{下一次研究的优先级}

第一优先级不是追加候选，而是冻结两项未命名类别的有限产品清单，并完成十项产品的当前估值和执行资料。第二优先级是对 ETF 建立久期/当前持仓和现金质量包，对 REIT 建立资产级运营、债务、资本开支、同日估值和实际分派桥。第三优先级才是独立回放与风险预算审查。届时每一只产品仍需单独通过，不存在“策略整体通过即自动放行”的机制。


\appendix
\chapter{名词解释}
\section{为何需要名词解释}

本附录定义本报告使用的产品与计算术语。定义用于避免把产品名称、披露口径或历史数据误读为承诺；不替代基金合同、招募说明书、交易规则或税务意见。

\begin{longtable}{L{3.0cm}L{11.5cm}}
\toprule
\textbf{术语} & \textbf{本报告中的含义} \\
\midrule
\endfirsthead
\toprule
\textbf{术语} & \textbf{本报告中的含义} \\
\midrule
\endhead
\textbf{ETF} & 交易型开放式指数基金。份额在交易所交易，基金可通过申购赎回机制运作。交易价格与基金净值可能不同；交易所成交不等于一定可按目标规模成交。\\
\textbf{公募 REIT} & 公开募集基础设施证券投资基金。本案中的 REIT 通过基金、专项计划和项目公司等结构持有基础设施相关权益，现金流来自收费公路、商业/住房租赁、产业园或发电等。它不是固定收益存款，也不是保证本金的产品。\\
\textbf{NAV（每份净资产）} & 按特定会计/估值日计算的基金每份净资产。NAV 的日期、估值方法和资产范围会影响可比性；历史 NAV 不能直接与不同日期的市场价格相减得出“当前折溢价”。\\
\textbf{IOPV} & 盘中估算净值。它通常用于辅助观察 ETF 交易价与估算资产价值的关系，但不替代经确认的基金 NAV，也不单独证明可执行性。\\
\textbf{价格/NAV 日期错配诊断} & 用某日市场价格与较早披露的 NAV 计算的算术差异。本报告仅把它当作提示进一步核验的诊断；不是当前溢价/折价、公允价值、目标价或交易信号。\\
\textbf{可供分配金额} & 公募 REIT 按法规、合同和披露口径计算的可分配金额。季度数值可能受到收现、费用、扩募、融资和时间分布影响，不能简单年化或等同于持有人未来现金。\\
\textbf{实际分派} & 已公告或已实施的现金分配金额。它仍需结合登记日、到账日、税费、份额、法律性质及是否存在返还本金等资料，才能进入“可复投净现金”统计。\\
\textbf{管理人年化率} & 发行人在某一披露期将历史期间数据按其定义年化得到的比率。本报告保留该披露事实，但不将其视为前瞻收益率、现金承诺或总回报预测。\\
\textbf{OCF（经营活动现金流）} & 财务报表中的经营活动产生的现金流量。它与净利润、可供分配金额和实际分派之间存在会计、营运资本、融资及资本开支差异，不能彼此替代。\\
\textbf{WALT} & 加权平均剩余租期。较短的 WALT 意味着更多租约较快进入续约/重租窗口；高出租率并不消除有效租金与续租风险。\\
\textbf{DV01 / 久期 / 凸性} & 债券或债券 ETF 对利率变化的敏感度指标。久期描述一阶价格敏感度，DV01 表示利率小幅变动的价值影响，凸性描述非线性部分。它们是评估 511010 的必要字段，而非由“国债”二字自动得到。\\
\textbf{TWRR} & 时间加权回报率。在外部现金流发生时划分子期间并链式连接，用于衡量策略本身回报。\\
\textbf{MWRR / XIRR} & 资金加权回报或带日期的内部回报率，反映特定投资者实际入金、出金和时间点；应与 TWRR 分开报告。\\
\textbf{滚动 12 个月净现金} & 本案对 3\% 目标的唯一候选口径：实际结算的合格现金，扣除税费、领取费用、持续成本和记录的分派下调，除以期初组合资本。卖出本金、未实现收益、外部资金和借款不计入。\\
\textbf{q=0} & 本报告的研究控制标记，表示该产品没有被批准为非零数量/权重，保持观察，不代表系统已经提交或撤回任何交易指令。\\
\textbf{同日流动性/容量} & 在同一交易日、以明确目标规模评估的买卖价差、盘口深度、成交可得性、冲击成本和市场做市情况。单日成交额本身不能替代这些字段。\\
\bottomrule
\end{longtable}

\sourcenote{定义结合产品法律文件、监管资料和本案计量控制。REIT 分派政策来源：\sourceid{S01, S12--S25}。}


\chapter{来源、审计轨迹与披露}
\section{可追溯来源包}

所有材料均已归档在本案 \texttt{sources/} 目录，派生核验、产品卡、市场和估值控制文件位于 \texttt{data/} 目录。文中以 \sourceid{Sxx} 标识来源族，而不把长网址塞入正文。完整 URL、本地相对路径、文档日期、预期用途、SHA-256 和限制在来源登记中；S27/S28 的文件级拆分见原子来源清单；每项关键主张到来源族和控制文件的映射在主张审计表中。

\begin{exhibitbox}[表 15：来源族与用途]
\small
\begin{tabularx}{\exhibitboxwidth}{L{2.0cm}L{4.0cm}X}
\toprule
\textbf{来源 ID} & \textbf{材料} & \textbf{本报告用途与限制} \\
\midrule
S01--S03 & CSRC/国务院/财政部公开政策材料 & REIT 分派制度及保护性流动性类别的边界；不识别具体流动性产品。\\
S04--S11 & 三只 ETF 的年报、Q2、产品资料概要/招募材料、指数方法 & 条款、费率、年报/Q2 NAV、历史表现、分配和已披露持仓；Q2 未审计且非当前持仓。\\
S12--S25 & 七只 REIT 的年报、Q2、法律与运营材料 & 项目资产、运营、现金流、历史分派、债务线索和更名/扩募；Q2 未审计，季度口径不可年化。\\
S26 & 508096 2026-02-13 更名公告 & 确认当前名称为中航京能国际能源REIT；不支持估值或准入。\\
S27 & 上交所 2026-08-12 同日快照 & 最后价、成交量/额、OHLC；不含点差、深度、IOPV/NAV、容量、税务或账户。\\
S28 & 全部产品的 PDS/招募材料归档 & 法律与产品条款交叉核验；不提供实时估值或实施批准。\\
\bottomrule
\end{tabularx}
\sourcenote{完整来源登记见 data/source\_registry.md/json；S27/S28 文件级来源清单见 data/legal\_document\_manifest\_20260812.md/json。}
\end{exhibitbox}

\section{卖方资料与失败探针}

本案还保留了 1 份原始券商 PDF、4 个产品级转载/预览和 3 个行业资料页面，以及失败检索记录。它们不构成可用的一致预期：没有一项为全部十个产品提供可审计的当前评级、目标、前瞻现金/收益预测和估值方法。因而每项卖方权重为零，报告没有伪造一致预期、目标价或隐含上行空间。详见本案的数据室中的卖方共识包及失败探针归档。

\section{尚未耗尽的实施资料清单}

尚未取得的资料不是被忽略，而是被记录为实施阻断项：同日 NAV/IOPV、REIT 资产估值桥、点差/深度/20 日容量、账户/税务/结算、原始卖方覆盖、完整历史回放数据，以及两项未命名类别的有限产品清单。每项都有下一步取证路径，详见 \texttt{source\_exhaustion\_log.md/json}。在这些字段补齐前，报告的正确落点是\textbf{证据备忘录 / 观察清单}，不是投资执行。

\begin{disclosurebox}[重要披露]
\small
\textbf{研究框架通过}仅指证据包可发布，不等于配置批准。所有市场产品均为 q=0；5\% 总回报与 3\% 可复投净现金目标\textbf{并不保证实现}。实施\textbf{尚未获批，亦未获授权}。
\end{disclosurebox}


\end{document}
